\addcontentsline{toc}{section}{\normalsize{ABSTRACT}}
\begin{center}
{\Large ABSTRACT}\\[24pt]
\textbf{SUBCHAT TREES: A SCALABLE HIERARCHICAL ARCHITECTURE FOR MULTI-THREADED DIALOGUE AND CONTEXT ISOLATION IN LLMS}
\end{center}
\justifying
Current LLM chat interfaces use linear transcripts to communicate with the user. In linear conversation every topic, sub-topic and even follow-up is appended to the same chronological history. This design is useful for short single-topic conversations but it becomes fragile in realistic long-term conversations. Where users switch between task/topic, return or mention earlier topic or ask nested follow-ups. Human conversations naturally follow a hierarchical structure but modern day AI interface lack conversational flow control in their interfaces. As a result, Unrelated context can leak across topics and older context can be evicted or buried. Current studies try to address this problem by increasing context length and models complexity. But Human conversation lacks its natural flow in chat interface. We see it as a design problem and propose Subchat Trees, a user-facing nested branching architecture for conversation management. Instead of treating the chat as one flat linear conversation, SubchatTree represents dialogue as a tree of parent and child subchats, where each node maintains its own local buffer, rolling summary and  a shared vector-indexed long-term memory. A user can create a focused subchat from a selected parent excerpt, continue a side discussion in that branch and later return to parent without contaminating sibling or ancestor contexts. This external architecture is model-agnostic and can imporve LLM performance without fine-tuning, as the model receives a cleaner branch-local context at inference time. The approach was evaluated on a controlled 353 turn Variable-Block Randomized Interleaving (VBRI) stress-test benchmark derived from five LMSYS-Chat-1M conversations, covering 59 topic labels and 43 topic-recall probes. Across Qwen2.5 7B, 14B and 32B models and buffer sizes $C \in \{5,10,20,40\}$. SubchatTrees improves context-isolation F1 over the linear baseline in every model-buffer setting. The best performing model was Qween32B at $C=20$, with an F1 score of 85.2\% and pollution rate was 16.2\%. Recall probes also show stronger topic retention. Average ROUGE-1 improves over baseline in all tested settings, including 0.2476 to 0.4187 for 14B at $C=20$ and 0.3338 to 0.4283 for 32B at $C=20$. A reduction in input token per correct answer is also observed across all tested configurations. The results suggest that complex conversations need explicit external structure, not only longer context windows.
